\chapter{Diaper Rash}

\section*{Definition}
Diaper dermatitis is an acute inflammatory eruption of the skin covered by a diaper, most commonly irritant contact dermatitis resulting from prolonged exposure to urine and feces. It affects up to 50\% of infants and toddlers at some point during the diaper-wearing period.

\section*{Pathophysiology}
Prolonged skin occlusion and moisture disrupt the stratum corneum barrier, increasing permeability and susceptibility to irritants. Fecal enzymes (urease, lipase, protease) generate ammonia and bile salts that raise skin pH, further damaging keratinocytes. Macerated skin becomes colonized by \textit{Candida albicans}, bacteria, or both, converting simple irritant dermatitis into secondary infection.

\section*{Etiology}
\begin{itemize}
  \item \textbf{Irritant:} Prolonged contact with urine and feces, infrequent diaper changes, diarrhea
  \item \textbf{Infectious:} \textit{Candida albicans} (satellite pustules), \textit{Staphylococcus aureus}, streptococcal perianal infection
  \item \textbf{Allergic/contact:} Fragrances, preservatives in wipes or diapers, topical preparations
  \item \textbf{Underlying conditions:} Seborrheic dermatitis, atopic dermatitis, psoriasis, zinc deficiency
\end{itemize}

\section*{Indications for Evaluation}
Seek care if:
\begin{itemize}
  \item Rash is severe, weeping, or accompanied by fever or irritability
  \item Satellite papules/pustules (suggests candidal superinfection)
  \item No improvement after 5--7 days of conservative management
  \item Poor feeding, dehydration, or signs of systemic illness
\end{itemize}

\section*{Related Symptoms}
\begin{itemize}
  \item Perianal erythema, papules, pustules, scaling, erosions, intertrigo
\end{itemize}
\textbf{Note:} Candida diaper dermatitis classically spares the skin folds, whereas irritant dermatitis involves the convex surfaces of the buttocks and thighs.

\section*{Self-Care / Home Management}
\begin{itemize}
  \item Change diapers every 2--3 hours and immediately after soiling; allow diaper-free time 10--15 minutes daily
  \item Cleanse gently with warm water and soft cloth; avoid commercial wipes with alcohol or fragrance
  \item Apply barrier ointment (zinc oxide 40\%, petrolatum) at each change
\end{itemize}

\section*{Diagnostic Approach}
\begin{itemize}
  \item Clinical diagnosis by inspection of diaper area distribution and morphology
  \item Potassium hydroxide (KOH) preparation of scale or pustule contents to identify pseudohyphae
  \item Bacterial culture if purulent drainage or suspected secondary infection
\end{itemize}

\section*{Treatment / Management Overview}
\begin{itemize}
  \item Irritant dermatitis: Barrier creams, frequent changes, low-potency topical steroid (hydrocortisone 1\% BID \ensuremath{\times}3--5 days)
  \item Candidal dermatitis: Topical antifungal (nystatin, clotrimazole, miconazole) BID--TID for 7--14 days
  \item Bacterial superinfection: Topical mupirocin or oral antibiotics based on culture
  \item Avoid combination antifungal--steroid products containing potent steroids (betamethasone) without prescription
\end{itemize}

\clearpage
